\documentclass[12pt]{article}
\usepackage[utf8]{inputenc}
\usepackage[T1]{fontenc}
\usepackage{hyperref}
\title{Qiubei Zhuang Language Resources for K-12 Curriculum Development}
\author{Compiled by Claude Sonnet 5.5 (claude-sonnet-5-5) for Andr\'as Kornai}
\date{}
\begin{document}
\maketitle

\section{Introduction}
Qiubei Zhuang (ISO 639-3 zqe, Glottocode qiub1238) is a Northern Tai language of Yunnan, China. It is part of the Yei Zhuang complex, and its speakers are also known as the Sha. Estimates of the number of speakers are 122,000 (Wikipedia) and 140,000 (Global Recordings Network). Speakers are concentrated in Qiubei, Shizong and Guangnan counties. The searches turned up reference entries and religious material, but no educational resources.

\section{Primary Educational Resources}
No primers, textbooks or school programmes for this variety were found in search results.

\section{Linguistic and Reference Materials}
\begin{itemize}
\item Wikipedia, Yei Zhuang language (covers the code zqe): \url{https://en.wikipedia.org/wiki/Yei_Zhuang_language}
\item A Zhuang-language web dictionary and writing resource (\emph{Zhuangwen}) turned up in the search. It concerns standard Zhuang and may not cover Qiubei specifically: \url{https://www.webonary.work/zhuangwen/overview/foreword/}
\end{itemize}

\section{Religious and Community Resources}
\begin{itemize}
\item Joshua Project profile, Qiubei Zhuang in China. It reports a Bible-portions translation (2024), a Jesus Film in the language, and no New Testament or complete Bible: \url{https://joshuaproject.net/people_groups/19806}
\end{itemize}

\section{Digital and Audio Resources}
\begin{itemize}
\item Global Recordings Network listing for Zhuang, Qiubei. Audio recordings are listed as unavailable at the time of the search: \url{https://www.globalrecordings.net/en/language/zqe}
\end{itemize}

\section{Curriculum Development Context}
The language is reported to be used by the whole community, but it is not known to be taught in schools. Standard Zhuang has its own educational and orthographic infrastructure, which a Qiubei curriculum might build on, but this was not verified in the searches. Dictionaries, grammars and learning materials specific to Qiubei Zhuang were not found, so they would need to be created.

\end{document}